\subsection*{章末確認問題}
\addcontentsline{toc}{subsection}{12 章末確認問題}
\begin{enumerate}
\item ソフトウェア構成管理が、単なる版管理ツールの利用ではない理由を説明せよ。
\item 構成品目とソフトウェア構成品目の違いを説明せよ。
\item 基準版とは何か。なぜ正式な変更統制なしに変更してはいけないのか。
\item 変更要求、ソフトウェア変更要求、構成管理委員会の関係を説明せよ。
\item 逸脱と免除の違いを説明せよ。
\item 構成状態記録・報告が、品質保証、セキュリティ、規制対応に役立つ理由を説明せよ。
\item 機能構成監査と物理構成監査の違いを説明せよ。
\item ソフトウェアビルドを再現可能にするために、どのような品目を構成管理下に置くべきか。
\item ソフトウェア部品表が、変更影響分析やセキュリティ対応に役立つ理由を説明せよ。
\item 小規模プロジェクトと全社的な開発組織では、構成管理ツールに求める範囲がどう変わるか。
\end{enumerate}
\begin{notebox}{解答の方向性}1 は「識別、変更統制、状態記録、監査、リリース管理を含む」ことを述べる。2 は「一つの管理単位」と「ソフトウェア関連の管理単位」の関係で説明する。3 は「承認済み構成を示す基点」であることを述べる。6 は「証跡を残す活動」であること、7 は「仕様どおりの機能か」と「実物と文書が一致するか」の違いで説明できる。\end{notebox}
